\subsection{Removing the negative corner when its residual class vanishes}
\label{sec:rank-three-corner-elimination}

We use the Hecke-track handoff in the actual fixed-inertia
rank-three presentation (DC3)--(DC7).  This is a refinement of
the point comparison, not a new variable-weight automorphy theorem.

\begin{proposition}[Integral negative-corner reduction]
\label{prop:rank-three-corner-elimination}
In the original normalized rank-three ring, suppose
\(H^1_L(-2)=0\), with exactly the local conditions used in
(DC4).  Then \(n_{-2}=0\), there are no equations of weight
\(-2\) in the list (DC6), and every two-sided point is type I,
in either coefficient characteristic.  This also applies
to two-sided points of the finite-sum graded algebra
of the ring.
\end{proposition}
\begin{proof}
A relative tangent of weight \(-2\) is a cocycle in the
negative corner line with the original local conditions.
It is a coboundary by hypothesis.  Normalization at complex
conjugation forces its primitive to be fixed by conjugation.
On the corner line the negative reversing transpose acts
as minus one.  Since the residue characteristic is odd,
the primitive, and hence the tangent, is zero.
The minimal relative presentation in (DC5) therefore
has no variable of weight \(-2\).
As the defect in that weight in (DC4) is zero, (DC6)
has no equation in that weight either.

Every homogeneous monomial of negative degree in the
minimal power series presentation contains a variable
of weight \(-1\), since all the other variables have
nonnegative weights.  The variables were chosen among
actual entry displacements.  In weight \(-1\) these are
the two adjacent lower entries.  The same divisibility
statement holds for a homogeneous series: partition its
monomials according to the first such variable dividing
them, then factor out that variable.  This is a finite
sum of homogeneous series times actual adjacent entries.
It proves the ideal inclusion both in the complete ring
and in the algebra of finite sums of homogeneous series,
without deleting any nilpotent.

By the marked projector classification in the preceding
integral count, a two-sided type-II point has its outer
two-dimensional irreducible block separate from the
middle line, so both adjacent lower entries vanish
identically.  The ideal inclusion forces all negative
weights to vanish, contradicting two-sidedness.
This proof uses neither rationalization nor a passage to
the reduced scheme.
\end{proof}

\paragraph{Explicit range in the RH frame.}
Write \(\psi=\psi_2\psi_{\rm odd}\), with
\(\psi_{\rm odd}^2=1\) and \(\psi_2\) primitive of conductor
\(2^b\), as in (DH1).  The residual negative corner
is \(\overline{\eta_4\epsilon^{-2e}\psi^{-2}}\).
It is odd and its primitive Teichm\"uller lift has odd
conductor \(f_0\mid\ell\); its dyadic conductor is
\(2^{b-1}\) for \(b\ge4\).
Take \(d=\operatorname{ord}_4(\ell)\) and
\(h=v_2(\ell^d-1)\).
The odd factor has order dividing \(\ell-1\).
If \(\ell\equiv1\bmod4\), then \(2^h\le\ell-1\);
otherwise \(2^h\le2(\ell+1)\), by the factorization
of \(\ell^2-1\).
In particular
\[
            2^b\ge8\ell(\ell+1)^2
       \quad\Longrightarrow\quad
            f_0\,4^h\le 2^{b-1}.
\]
Apply Corollary~\ref{cor:selmer-dyadic-twist-cocycle}.
Away from \(\ell\) the original inertia kernels have
prime-to-\(\ell\) index, and at \(\ell\) the local
coboundary vanishes on the kernel of the coefficient
character.  Thus its literal relative-index hypotheses
hold.  This proves \(H^1_L(-2)=0\) in the displayed range.
At \(\ell=3\) the earlier proof of \texttt{Xpp}
in Theorem~\ref{thm:rh-dyadic-host-frame} already gives
the conclusion under (DH1), without the stronger
threshold.  In these ranges the abstract length comparison
needs only the irreducible point spectrum and its rigidity,
not the type-II endpoint vanishing.  For the intermediate
conductors retain the full preceding proof.

\paragraph{The Dirichlet application.}
For any admitted parent fixed-inertia rank-three frame,
let the negative corner be
\(\bar\gamma=\overline{\eta\Psi^{-2}}\), and suppose
its primitive finite-order lift is odd, with dyadic
conductor \(2^{b-1}\), odd conductor \(f_0\), and
prime-to-\(\ell\) order.  For any \(d\ge1\) such that
\(4\mid\ell^d-1\) and the odd factor's order divides
\(\ell^d-1\), the sufficient criterion is
\[
                      f_0\,4^{v_2(\ell^d-1)}\le 2^{b-1}.
\]
The same original local relative-index checks are
necessary.  If the moving odd factors are quadratic,
squaring removes them and \(f_0\mid M_0\ell\) with
\(M_0\) fixed (including the fixed multiplier's odd
conductor).  When the odd factor's order divides
\((\ell-1)s\) for a fixed \(s\) prime to \(\ell\), one
may choose \(d=2s\operatorname{ord}_s(\ell)\), with
\(\operatorname{ord}_1(\ell)=1\).  In fact the geometric
sum for this \(d\) is divisible by \(s\), by grouping
it into blocks of length \(\operatorname{ord}_s(\ell)\).
Also \(d\) is even, so the elementary 2-adic valuation
formula gives \(2^{v_2(\ell^d-1)}\ll_s\ell+1\).
This yields a sufficient threshold
\(C_{M_0,s}\ell(\ell+1)^2\).
The exact displayed criterion is more useful than this
coarse sufficient one at individual primes.

This reduction does not assert that every Dirichlet datum
admits the fixed-inertia frame, nor that the auxiliary
variable-weight ring has no reducible components.  It
does not change the rank-four construction or the already
assembled ordinary-RH weighted length estimate.
The strong-irreducibility test in the same handoff is
available only under its distinct-ratio and nonscalar
open-subgroup hypotheses.
